\chapter{Ứng dụng hệ thống gợi ý nhạc theo mood}
\label{ch:ungdung}

\section{Tổng quan}

Agent V1 chạy vòng lặp suy luận bốn bước: phân tích yêu cầu, gọi công cụ, quan sát kết quả, trả lời người dùng (LangGraph). API REST phục vụ chat; server MCP phục vụ công cụ tìm kiếm nội bộ.

Corpus retrieval: 2.261 bài kèm ma trận embedding 768 chiều.

\section{Chính sách mood}

Agent phân biệt \textbf{mood hiện tại} (cảm xúc người dùng mô tả) và \textbf{mood mục tiêu} (mood nhạc cần gợi ý). Sáu mood chuẩn trùng tập crawl.

\section{Công thức chấm điểm retrieval}

\begin{table}[htbp]
  \centering
  \caption{Trọng số scoring retrieval V1}
  \label{tab:scoring}
  \begin{tabular}{@{}ll@{}}
    \toprule
    \textbf{Thành phần} & \textbf{Trọng số} \\
    \midrule
    Tương đồng ngữ nghĩa (cosine embedding) & 0{,}55 \\
    Khớp mood & 0{,}20 \\
    Khớp thể loại & 0{,}20 \\
    Khớp tag & 0{,}05 \\
    \bottomrule
  \end{tabular}
\end{table}

Nghệ sĩ lọc khớp chính xác, không embed. V1: trường thể loại rỗng trên corpus.

\section{API và kiểm thử}

API chat trả câu trả lời tự nhiên, danh sách gợi ý và lịch sử gọi công cụ. Mỗi request ghi log trace agent.

Bộ kiểm thử: 17 module pytest, chạy offline không cần khóa API.

\section{Liên kết dữ liệu và agent}

\begin{table}[htbp]
  \centering
  \caption{Dữ liệu pipeline và vai trò trong agent}
  \label{tab:data-agent}
  \begin{tabularx}{\textwidth}{@{}L{4cm}X@{}}
    \toprule
    \textbf{Dữ liệu} & \textbf{Vai trò} \\
    \midrule
    Tóm tắt metadata bài hát & Text embed cho retrieval \\
    Nhãn mood & Lọc và chấm điểm (0{,}20) \\
    Tag & Chấm điểm bổ sung (0{,}05) \\
    Ma trận embedding & Cosine similarity (0{,}55) \\
    URL preview Deezer & Phát thử trên giao diện \\
    Tập mẫu 30 bài & Kiểm thử offline \\
    \bottomrule
  \end{tabularx}
\end{table}

Đặc trưng âm thanh 54 chiều có trong tập huấn luyện nhưng chưa dùng trong scoring V1.
